\begin{ejercicio}
    ¿Cuáles son los límites de integración de la región $W$ para las ecuaciones 3, 4 y 5 (la región es un cono simétrico con tapa)?

\begin{enumerate}
    \item[(a)] Describan cada paso del proceso a través del cual fueron definiendo los límites de integración.
    \item[(b)] Incluyan todas las gráficas que consideren necesarias para ilustrar los límites de integración para cada sistema de coordenadas.
    \item[(c)] Argumenten de qué depende que sea más cómoda o preferible la resolución de una integral en unas coordenadas o en otras.
\end{enumerate}

\begin{align}
    &\int_{?}^{?} \int_{?}^{?} \int_{?}^{?} f(r, \theta, z)\, r\, dz\, dr\, d\theta \\
    &\int_{?}^{?} \int_{?}^{?} \int_{?}^{?} g(\rho, \varphi, \theta)\, \rho^2 \sin\varphi\, d\rho\, d\varphi\, d\theta \\
    &\int_{?}^{?} \int_{?}^{?} \int_{?}^{?} h(x, y, z)\, dx\, dy\, dz
\end{align}
\end{ejercicio}

\begin{respuesta}
La región \(W\) se encuentra delimitada por la superficie

\[
f(x,y)=2e^{-(x-1)^2-y^2}
\]

y el disco

\[
x^2+y^2\leq 4.
\]

\subsection*{(a) Descripción de los contornos}

El problema nos dice que tomemos de eejmplo esto:

\[
f(x,y)=1.
\]

Ahora pasamos a sustituir la funcion que tenemos y nos queda:

\[
2e^{-(x-1)^2-y^2}=1.
\]

Ahora dividimos entre \(2\):

\[
e^{-(x-1)^2-y^2}=\frac{1}{2}.
\]

Despues aplicamos el logaritmo natural, la propiedad de la función inversa entre la exponencial natural y el logaritmo natural cancela las operaciones y nos permite quedarnos con solo los exponentes, entonces nos queda: 

\[
-(x-1)^2-y^2=\ln\left(\frac{1}{2}\right).
\]

Por otra parte tenemos que: 

\[
\ln\left(\frac{1}{2}\right)=-\ln 2,
\]

Sustituyengo y multiplicando por -1 ambos lados obtenemos: 

\[
(x-1)^2+y^2=\ln 2.
\]

Recordando que la ecuación general de un círculo es

\[
(x-h)^2+(y-k)^2=r^2,
\]

podemos escribir la ecuación obtenida como

\[
(x-1)^2+(y-0)^2=\ln 2.
\]

Al comparar ambas ecuaciones, tenemos

\[
h=1,\qquad k=0.
\]

Por lo tanto, el centro del contorno es

\[
{(1,0)}
\]

y su radio es

\[
{\sqrt{\ln 2}}.
\]


Por otra parte, si tenemos \(f(x,y)=c\) entonces los contornos están dados por

\[
(x-1)^2+y^2=\ln\left(\frac{2}{c}\right).
\]
Entonces podemos decir que todos los contornos estarán centrados en (1,0), y para que el contorno exista debemos tener que \(0<c\leq 2\). Si \(c=2\), el contorno se reduce al punto \((1,0)\), que es el máximo de \(f\).Si \(c>2\) no hay contorno.

\subsection*{(b) Área de la sección correspondiente al plano \(x=1\)}

Tenemos el plano: 

\[
x=1.
\]

Entonces ahora lo sustituimos en (x=1) en la función y nos queda: 

\[
f(1,y)=2e^{-(1-1)^2-y^2}
=2e^{-y^2}.
\]

Para obtener los limites de \(y\), utilizamos el disco: 

\[
x^2+y^2\leq4.
\]

Ahora volvemos a sustituir a x=1 y tenemos: 

\[
1+y^2\leq4,
\]

Estonces: 

\[
y^2\leq3.
\]

De manera que nos queda: 

\[
-\sqrt{3}\leq y\leq\sqrt{3}.
\]
Entonces la region W se encuentra entre el plano (z=0) y la superficie (z=f(x,y)), sobre el disco (x^2+y^2\leq4).
Cuando cortamos a W con el plano (x=1), la sección es la región del plano (x=1) comprendida entre (z=0) y la curva (z=2e^{-y^2}), con

\[
-\sqrt{3}\leq y\leq\sqrt{3}.
\]

Por lo tanto, el área de esta sección es:

\[
A=\int_{-\sqrt{3}}^{\sqrt{3}}2e^{-y^2},dy.
\]
\subsection*{(c) Escriban la integral iterada doble que de el volumen de W}

Tenemos lo siguiente: 
\[
f(x,y)=2e^{-(x-1)^2-y^2}
\]
y la superficie está sobre la región
\[
x^2+y^2\leq 4.
\]
Esto representaria un circulo de radio de 2.
Como \(f(x,y)>0\), el sólido \(W\) queda entre el plano \(z=0\) y la superficie
\(z=f(x,y)\), entonces podemos obtener el volumen mediante la integral doble:
Por lo tanto,
\[
V=\iint_R f(x,y)\,dA.
\]
En donde R seria el disco definido por:
\[
x^2+y^2\leq4.
\]
Entonces obtenemos limites de integracion para esto despejamos \(x\):
\[
x^2\leq4-y^2
\]
\[
-\sqrt{4-y^2}\leq x\leq\sqrt{4-y^2}.
\]
Además, como el círculo tiene radio \(2\):
\[
-2\leq y\leq2.
\]
Por lo tanto, la integral iterada que representa el volumen de (W) es:
\[
{
V=
\int_{-2}^{2}
\int_{-\sqrt{4-y^2}}^{\sqrt{4-y^2}}
2e^{-(x-1)^2-y^2}
\,dx\,dy
}
\]
    \textbf{Integrantes involucrados: Alan Blancas Ochoa }
\end{respuesta}
